\subsection*{Teil C: Knobelaufgaben (25 Minuten)}

\begin{enumerate}[label=\textbf{C\arabic*}]

    \item Ein Würfel und eine Münze werden gleichzeitig geworfen. Zeichne ein Baumdiagramm (erst Würfel, dann Münze) und gib an, wie viele Ergebnisse es gibt.

    \vspace{5cm}

    \item Ein Baumdiagramm hat auf der ersten Stufe drei Äste mit den Beschriftungen 1, 2 und 3. Von jedem dieser Äste gehen auf der zweiten Stufe zwei Äste mit den Beschriftungen W und Z ab. Erfinde ein passendes Zufallsexperiment.

    \vspace{0.3cm}
    \antwortlinie{\linewidth}\\[0.4cm]
    \antwortlinie{\linewidth}

    \vspace{0.6cm}

    \item Eine Münze wird viermal geworfen. Wie viele Ergebnisse gibt es? Finde die Antwort, ohne das ganze Baumdiagramm zu zeichnen, und erkläre deinen Rechenweg.

    \vspace{0.3cm}
    \antwortlinie{\linewidth}\\[0.4cm]
    \antwortlinie{\linewidth}

    \vspace{0.6cm}

    \item Vergleiche B3 und B4. Warum hat das Baumdiagramm in B4 weniger Pfade? Erkläre in einem Satz.

    \vspace{0.3cm}
    \antwortlinie{\linewidth}\\[0.4cm]
    \antwortlinie{\linewidth}

    \vspace{0.6cm}

    \item Max sagt: „Ob mein Lieblingsverein am Samstag gewinnt, ist ein Zufallsexperiment, denn das Ergebnis kann man nicht vorhersagen.“ Hat Max recht? Prüfe alle vier Bedingungen.

    \vspace{0.3cm}
    \antwortlinie{\linewidth}\\[0.4cm]
    \antwortlinie{\linewidth}\\[0.4cm]
    \antwortlinie{\linewidth}

\end{enumerate}
